\begin{textsample}{POS Dim 1 – vem\_ed\_21 – Score 7.00 – vem\_ed\_21/t107.txt}  \label{ex:f1_pos_167}
Comunicações Do quebra-gelo à institucionalização Regiane Souza Camargo Moreira apresentou o trabalho “¿ Cómo evaluar el \textbf{impacto} de un rompe-hielo en proyecto COIL?” (Como avaliar o \textbf{impacto} de um quebra-gelo no projeto COIL?), elaborado em coautoria com Javier Guerrero e Sandra Cruz (Uniminuto / Colômbia).
Destacou \textbf{que} o quebra-gelo \textbf{promove} \textbf{conexão} e \textbf{diálogo} intercultural entre os participantes, daí a \textbf{importância} de seu planejamento. \textbf{É} \textbf{essencial} se preparar para poder se expor, especialmente nos aspectos pessoais e culturais. “O quebra-gelo não garante o sucesso do projeto, mas ajuda a reforçar o trabalho em equipe e a evolução do COIL”, afirmou Regiane, \textbf{que} \textbf{compartilhou} pesquisa realizada com 30 professores e 4 coordenadores COIL.
Ana Carolina Freschi fez duas comunicações relacionadas à institucionalização dos projetos COIL.
A primeira \textbf{foi} “GSL Classroom and the Chamber of Secrets: The Journey throughout Institutionalizations of COIL Initiatives” (Intercâmbios Virtuais e a Câmara Secreta: a jornada pelas institucionalizações das iniciativas COIL) com Ricardo Lyle Bañuelos, Gisselle Morales e Monserrat Balandrano (Tecnológico de Monterrey / México), Dan Nolan (University of Minnesota System / EUA) e Natalia Dyba (University of Washington Bothell / EUA).
Os apresentadores \textbf{destacaram} a \textbf{importância} da melhoria \textbf{contínua} dos \textbf{processos} de institucionalização, com diretrizes e requisitos detalhados.
O Tecnológico de Monterrey, por exemplo, trabalha com chamada de \textbf{propostas} para conhecer os planos de colaboração internacional dos professores interessados em iniciativas GSL (Global Shared Learning, \textbf{que} \textbf{é} como a instituição chama os Intercâmbios Virtuais).

% matched lemmas: compartilhar, conexão, contínuo, destacar, diálogo, essencial, impacto, importância, processo, promover, proposta, que, ser
\end{textsample}
